\documentclass[11pt]{article}
\usepackage[margin=2.3cm]{geometry}
\usepackage{amsmath,amssymb,booktabs}
\allowdisplaybreaks
\input{../build/R02_numbers.tex}
\title{XVI-R R02: computed BRST-trivial span, classification of the 43 unclassified directions of $N_g$,\\
and a canonical candidate basis of invariant evanescent operators at $\kappa^3$}
\author{XVI-R (local landmark; imports XVI-Q frozen at Q11; not published)}
\date{}
\begin{document}\maketitle
\begin{abstract}
At order $\kappa^3$, three points, pole level, the evanescent defect space $N_g$ ($\dim 104$) left by XVI-Q is examined
with a \emph{computed} $\sigma$-map on the full monomial space of local ghost-number $-1$ functionals used since Q08,
rather than with individually tested generators. Of the 43 directions that Q11 left unclassified, \NT{} lie in the computed
trivial span; none is blind; \NHC{} remain. The \NHC{} remaining directions lie outside the linear span of
\emph{all} computed $\sigma$-images and blind coordinates (\NR{}/\NHC{} robust), and we give a canonical basis with explicit
tensor structures. Because the $h\cdot h\cdot c$ $\sigma$-composition gap of Q11 remains open, these are
\textbf{cohomology candidates}, not proven nontrivial classes; no direction is certified physical. Outcome: \textbf{Partial}.
\end{abstract}

\section{Scope and premises}
The only premise is XVI-Q frozen at Q11, imported by its source-only zip (sha \texttt{1888f327\ldots}) and rebuilt inside
this build. XVI-N, XVI-O, XVI-P and XVI-Q are read-only. OpenAI.tex and the newer companion sheets are not premises.
O.1 remains CLOSED only in its scoped sense (NOT\_CLOSED elsewhere), bulk symmetry is not restored, and the insertion stays
frozen. Missing is never treated as zero. One disclosed minimal revision (R02-R1): the R01 gate's unpinned-file check now
excludes files introduced by R02, which the R02 gate pins.

\section{Set-up}
The unknowns are the pole-level evanescent coefficients $X=(X_{K hc},X_{\bar c ch},X_{c^*cc},X_{hhh})$ in a fixed monomial
basis (821 coordinates; every coordinate carries at least one hatted, $(d-4)$-dimensional contraction). The linear system
$A$ (rows $ccK$, $cc\bar c$, $hhc$, plus Bose symmetry of $X_{hhh}$) and the ghost equation $A_g$
($\Gamma_{\bar c\ldots}=\rho\,\Gamma_{K\ldots}[K\to T(\bar c)]$, $\rho=-2$) define
$N_g=\ker\binom{A}{A_g}$, $\dim N_g=104$. Q11 found the decomposition
$104 = (\text{individually passing }\sigma\text{-images}) + (\text{blind}) - (\text{overlap}) + 43$, with $\dim W_{\rm Q11}=\WQ$.

\paragraph{The computed $\sigma$-map.} The ghost-number $-1$ local functionals are spanned by
\begin{itemize}
\item $F$: $\tilde K\cdot F(h,h)$, momentum degree 2, with $\tilde K=K+\bar c F'(h)$, so its $\bar c$ partner is $\rho\,T(\cdot)$;
\item $G$: $c^*\cdot G(c,h)$, momentum degree 2 (with the $c^*cc$ piece from $h\to s_0h$);
\item $P$: $\bar c\cdot\Phi(h,h)$, momentum degree 3 (standalone; the ghost equation is imposed on the image, not assumed).
\end{itemize}
Each generator $Y$ is mapped to $\sigma(Y)=$ (antifield-type pieces) $-$ (field-type pieces), with the channel operators
($D$: $K\to\delta S_2/\delta h$; $A$: $h\to s_0h$; $E$: $\bar c\to\sigma_F F(h)$) and signs derived in Q09/Q11. No
normalisation is fitted. The result is a matrix $\Sigma$ (821 $\times$ \NY). No generator leaks outside the basis.

\section{Results}
\subsection{Computed trivial span}
Let $Z=\ker(\binom{A}{A_g}\Sigma)$ and $T=\Sigma Z$. Then $T$ is the set of $\sigma$-images, including arbitrary linear
combinations of generators that fail individually, that satisfy every constraint. We find
$\operatorname{rank}\Sigma=\RS$, $\dim Z=\ZY$ and $\dim T=\DT$. If the ghost equation is dropped, $\dim(\Sigma\ker(A\Sigma))=\DTA$.
\input{../build/R02_tab.tex}
The standalone $P$ family contributes nothing to $T\cap N_g$ on its own. This is consistent with the ghost equation, which
admits $\bar c$ only through $\tilde K$.

\subsection{Classification of the 43}
Set $W'=W_{\rm Q11}+T+\text{blind}_g$. We find $\dim(W'\cap N_g)=\WP$, so
\[ 43 \;=\; \underbrace{\NT}_{\text{trivial (computed }\sigma)} \;+\; \underbrace{0}_{\text{blind}} \;+\;
\underbrace{\NHC}_{\text{cohomology candidates}},\qquad \text{physical (proven)}=0 . \]
The \NT{} newly trivial directions are combinations of $F$ and $G$ images that fail one by one but whose combination
satisfies all constraints. This answers Q11's question of whether the 43 contain hidden trivial directions: \NT{} of them are.
\emph{Robustness:} $\dim\big(N_g\cap(\operatorname{span}\Sigma+\text{blind}_g)\big)=\WR=\dim(W'\cap N_g)$. No
combination of computed images, even one ignoring the constraints, reaches any of the \NHC{} candidates. Any class
assignment of individual directions depends on the choice of basis; the statement that is independent of that choice is the
dimension count above, together with the canonical basis below.

\subsection{Why the candidates are not certified}
A candidate would be a nontrivial cohomology class if (i) $\Sigma$ were the complete image of the Slavnov--Taylor
linearisation on local gh $-1$ functionals at this order, and (ii) $\sigma(Y)\subset N$ for every $Y$. Neither is proven.
(ii) fails as computed: the Q11 composition gap (individually, 123 of 194 $F$ images miss $N$, all in the $hhc$ row) is open.
R01 localised it: the defects are correctable only with mixed-sector, mostly hat-degree-1 corrections. A correct
$\sigma$ could therefore have images that reach the candidates. The candidate basis is thus an \emph{upper bound}
($\le\NHC$) on the invariant evanescent operators at this order, and it is conditional on the monomial Y-space being
complete.

\section{Canonical candidate basis}
Representatives are the reduced row-echelon form of $N_g$ modulo $W'$ in the fixed coordinate order (sectors $K,C,S,H$;
Q08 monomial order). The resulting basis is unique and does not depend on the choice of kernel basis.
Notation: $[K]$, $[C]$, $[H]$ label the amplitudes $X_{Khc}$, $X_{\bar c ch}$, $X_{hhh}$; $\epsilon_i$ are the graviton
polarisation (symmetric-tensor legs, polarised), $c_1$ the ghost, $\hat x\cdot\hat y$ the evanescent contraction, $p_1,p_2$ the independent momenta.
Six candidates involve the antifield $K$ (three purely $K$-sector, three mixing $K$, $\bar c$ and $hhh$); sixteen are pure $hhh$
evanescent structures.
\input{../build/R02_basis.tex}

\section{Open items}
\begin{enumerate}
\item Close the $h\cdot h\cdot c$ $\sigma$-composition gap. Until then, no candidate can be proven nontrivial and none is called physical.
\item Prove that the F/G/P monomial Y-space is complete at this order. Higher momentum degree and other index structures have not been included.
\item p22 $\kappa^2$ evanescent-CT consistency: OPEN, not attempted.
\item O.1 stays scoped; bulk symmetry is not restored; Route B is not discharged; XV-B and XIV-J remain unopened.
\end{enumerate}
\end{document}
